\documentclass[12pt]{article}
\usepackage[margin=1in, top=0.85in, bottom=0.85in]{geometry}
\usepackage{hyperref}
\hypersetup{
  colorlinks   = true,
  urlcolor     = blue,
  linkcolor    = black,
  citecolor    = black
}
\usepackage[skip=7pt]{parskip}

\begin{document}

\begin{flushright}
September 2, 2026
\end{flushright}

\vspace{0.5em}

\textbf{Subject: Submission to \emph{The Journal of Real Estate Finance and Economics}}

\vspace{0.5em}

Dear Editors,

\vspace{0.5em}

Please consider my manuscript, ``Keeping the House in the Family: Non-Market Transfers and the Supply of American Homes,'' for publication in \emph{The Journal of Real Estate Finance and Economics} as a Research Article.

The paper uses the near-universe of recorded U.S. residential deeds to document a channel of housing turnover that runs alongside, and interacts with, the market channel: transfers of homes between family members. Next, it leverages a natural experiment to show that a specific real-estate tax policy -- California's rules allowing an heir to inherit a parent's below-market property-tax assessment along with the house -- causally shifts homes between market-based and lineage-based allocation. This connects directly to the journal's core interests in housing supply, property taxation, and the frictions that govern how homes change hands.

The paper's main contributions are threefold. First, using 144 million residential deed records from 2007--2023, it provides the first national, deed-based accounting of intra-family housing transfers, separating arm's-length sales from family transfers, trust retitling, and same-person co-owner changes. Families moved nearly 20 million homes between relatives over this period, one for every 3.6 market sales. Second, linking parcels across owners shows family-transferred homes are unusually sticky: 62 percent of same-surname transfers have not resold a decade later, exceeding the retention rate of owner-occupant and investor purchases combined. Third, it exploits Proposition 19 -- which abruptly ended inherited-assessment portability for non-occupant heirs in February 2021 -- as a natural experiment: a synthetic control and monthly event study show family transfers surged 31 percent ahead of the deadline and then settled 26--33 percent below it, a missing flow of roughly 58,000 deeds per year matching the state's own estimate of tax-favored inherited properties. A simple keep-or-sell model explains why these homes have not yet resurfaced as estate sales.

Replication code reproducing every table and figure at the aggregated level will be made available in a public repository. 
The underlying deed and property-characteristics microdata are proprietary, obtained under a restricted academic license from CoreLogic (now Cotality).

I confirm that this manuscript is original, is not under review elsewhere, and that I have no competing interests to disclose.

Thank you for your consideration. I look forward to hearing from you.

Sincerely,\\[1ex]
Saani Rawat\\
Assistant Professor of Finance \& Real Estate\\
Marquette University\\
\href{mailto:saannidhya.rawat@marquette.edu}{saannidhya.rawat@marquette.edu}

\end{document}
